
\documentclass{article}
\usepackage{colortbl}
\usepackage{makecell}
\usepackage{multirow}
\usepackage{supertabular}

\begin{document}

\newcounter{utterance}

\centering \large Interaction Transcript for game `function\_detective', experiment `logic\_passive\_labeled\_pairs', episode 201 with mistral{-}large{-}3{-}azure.
\vspace{24pt}

{ \footnotesize  \setcounter{utterance}{1}
\setlength{\tabcolsep}{0pt}
\begin{supertabular}{c@{$\;$}|p{.15\linewidth}@{}p{.15\linewidth}p{.15\linewidth}p{.15\linewidth}p{.15\linewidth}p{.15\linewidth}}
    \# & \multicolumn{2}{c}{Player} && \multicolumn{2}{c}{Game Master} \\
    \hline

    \theutterance \stepcounter{utterance}  
    & & & \multicolumn{4}{p{0.6\linewidth}}{
        \cellcolor[rgb]{0.9,0.9,0.9}{
            \makecell[{{p{\linewidth}}}]{
                \texttt{\tiny{[P1$\langle$GM]}}
                \texttt{I have a hidden Python function for you to figure out.} \\
\\ 
\texttt{Your job is to infer what it does and then submit a correct Python implementation.} \\
\\ 
\texttt{Function signature:} \\
\texttt{(a: bool, b: bool) {-}> bool} \\
\\ 
\texttt{You have 15 turns total.} \\
\\ 
\texttt{I reveal candidate input{-}output pairs one at a time and tell you whether each one is correct.} \\
\\ 
\texttt{You must make your best SOLVE attempt no later than the second{-}to{-}last turn.} \\
\texttt{Do not wait until the very end.} \\
\\ 
\texttt{You can use these actions:} \\
\texttt{{-} JSON format described at the end of the prompt} \\
\texttt{{-} SOLVE: ```python ... ```} \\
\\ 
\texttt{Do not add import statements in your guess.} \\
\texttt{Anything you might need is already available, so just write the function body in hypothesis and return the answer.} \\
\\ 
\texttt{Preloaded examples:} \\
\texttt{1. (((True, True), True), True)} \\
\\ 
\texttt{Scoring:} \\
\texttt{{-} correctness is required for any non{-}zero score} \\
\texttt{{-} in interactive modes, solving with fewer information requests is better} \\
\\ 
\texttt{Input formatting reminders:} \\
\texttt{{-} strings must be quoted} \\
\texttt{{-} booleans must be True/False} \\
\texttt{{-} use commas between multiple values} \\
\texttt{{-} do not write any extra action tags in the same message} \\
\\ 
\texttt{Format example:} \\
\texttt{You write:} \\
\texttt{\{'current\_mode': <probing/confirming/don't know; choose whichever best describes your current state in the task>, 'input\_rationale': <brief description of why you chose this hypothesis>, 'hypothesis': <either the exact string "unknown" if you have no guess yet OR your current best guess as runnable Python matching the given signature, with a one{-}line docstring stating what it does. Never write placeholder code, if you cannot commit to a guess, use "unknown". Match the given signature.} \\
\texttt{e.g.} \\
\texttt{def f(x):} \\
\texttt{	"""Returns the absolute value of x."""} \\
\texttt{	return x if x >= 0 else {-}x} \\
\\ 
\texttt{Use "unknown" if you have no guess.>, 'confidence\_score': <score between 1{-}100. 100 means highest confidence in your hypothesis>\}} \\
\\ 
\texttt{Output only one valid json per turn.} \\
            }
        }
    }
     \\ \\

    \theutterance \stepcounter{utterance}  
    & \multicolumn{4}{p{0.6\linewidth}}{
        \cellcolor[rgb]{0.561,0.737,0.561}{
            \makecell[{{p{\linewidth}}}]{
                \texttt{\tiny{[P1$\rangle$GM]}}
                \texttt{```json} \\
\texttt{\{} \\
\texttt{    "current\_mode": "probing",} \\
\texttt{    "input\_rationale": "The preloaded example ((True, True), True) suggests that when both inputs are True, the output is True. This could be consistent with several logical operations (AND, OR, XOR, etc.), but AND is the most straightforward starting hypothesis. I will test this hypothesis with a case where one input is False.",} \\
\texttt{    "hypothesis": "def f(a: bool, b: bool) {-}> bool:\textbackslash{}n    \textbackslash{}"\textbackslash{}"\textbackslash{}"Returns the logical AND of a and b.\textbackslash{}"\textbackslash{}"\textbackslash{}"\textbackslash{}n    return a and b",} \\
\texttt{    "confidence\_score": 70} \\
\texttt{\}} \\
\texttt{```} \\
            }
        }
    }
    & & \\ \\

    \theutterance \stepcounter{utterance}  
    & & & \multicolumn{4}{p{0.6\linewidth}}{
        \cellcolor[rgb]{0.9,0.9,0.9}{
            \makecell[{{p{\linewidth}}}]{
                \texttt{\tiny{[P1$\langle$GM]}}
                \texttt{OUTPUT: (((False, False), False), False)} \\
            }
        }
    }
     \\ \\

    \theutterance \stepcounter{utterance}  
    & \multicolumn{4}{p{0.6\linewidth}}{
        \cellcolor[rgb]{0.561,0.737,0.561}{
            \makecell[{{p{\linewidth}}}]{
                \texttt{\tiny{[P1$\rangle$GM]}}
                \texttt{```json} \\
\texttt{\{} \\
\texttt{    "current\_mode": "confirming",} \\
\texttt{    "input\_rationale": "The new example ((False, False), False) is consistent with the logical AND hypothesis, as False AND False should return False. To further confirm, I will test a case where one input is True and the other is False, which should return False for AND.",} \\
\texttt{    "hypothesis": "def f(a: bool, b: bool) {-}> bool:\textbackslash{}n    \textbackslash{}"\textbackslash{}"\textbackslash{}"Returns the logical AND of a and b.\textbackslash{}"\textbackslash{}"\textbackslash{}"\textbackslash{}n    return a and b",} \\
\texttt{    "confidence\_score": 90} \\
\texttt{\}} \\
\texttt{```} \\
            }
        }
    }
    & & \\ \\

    \theutterance \stepcounter{utterance}  
    & & & \multicolumn{4}{p{0.6\linewidth}}{
        \cellcolor[rgb]{0.9,0.9,0.9}{
            \makecell[{{p{\linewidth}}}]{
                \texttt{\tiny{[P1$\langle$GM]}}
                \texttt{OUTPUT: (((True, False), False), True)} \\
            }
        }
    }
     \\ \\

    \theutterance \stepcounter{utterance}  
    & \multicolumn{4}{p{0.6\linewidth}}{
        \cellcolor[rgb]{0.561,0.737,0.561}{
            \makecell[{{p{\linewidth}}}]{
                \texttt{\tiny{[P1$\rangle$GM]}}
                \texttt{```json} \\
\texttt{\{} \\
\texttt{    "current\_mode": "confirming",} \\
\texttt{    "input\_rationale": "The latest example ((True, False), False) confirms the logical AND hypothesis, as True AND False returns False. This aligns perfectly with all observed examples so far. To be thorough, I will now test the remaining case ((False, True), False) to ensure symmetry.",} \\
\texttt{    "hypothesis": "def f(a: bool, b: bool) {-}> bool:\textbackslash{}n    \textbackslash{}"\textbackslash{}"\textbackslash{}"Returns the logical AND of a and b.\textbackslash{}"\textbackslash{}"\textbackslash{}"\textbackslash{}n    return a and b",} \\
\texttt{    "confidence\_score": 99} \\
\texttt{\}} \\
\texttt{```} \\
            }
        }
    }
    & & \\ \\

    \theutterance \stepcounter{utterance}  
    & & & \multicolumn{4}{p{0.6\linewidth}}{
        \cellcolor[rgb]{0.9,0.9,0.9}{
            \makecell[{{p{\linewidth}}}]{
                \texttt{\tiny{[P1$\langle$GM]}}
                \texttt{OUTPUT: (((False, True), False), False)} \\
            }
        }
    }
     \\ \\

    \theutterance \stepcounter{utterance}  
    & \multicolumn{4}{p{0.6\linewidth}}{
        \cellcolor[rgb]{0.561,0.737,0.561}{
            \makecell[{{p{\linewidth}}}]{
                \texttt{\tiny{[P1$\rangle$GM]}}
                \texttt{```json} \\
\texttt{\{} \\
\texttt{    "current\_mode": "confirming",} \\
\texttt{    "input\_rationale": "The final example ((False, True), False) confirms the logical AND hypothesis, as False AND True returns False. All four possible input combinations now align with the behavior of the AND operator. I am confident in this solution and will submit it.",} \\
\texttt{    "hypothesis": "SOLVE: ```python\textbackslash{}ndef f(a: bool, b: bool) {-}> bool:\textbackslash{}n    \textbackslash{}"\textbackslash{}"\textbackslash{}"Returns the logical AND of a and b.\textbackslash{}"\textbackslash{}"\textbackslash{}"\textbackslash{}n    return a and b\textbackslash{}n```",} \\
\texttt{    "confidence\_score": 100} \\
\texttt{\}} \\
\texttt{```} \\
            }
        }
    }
    & & \\ \\

    \theutterance \stepcounter{utterance}  
    & & & \multicolumn{4}{p{0.6\linewidth}}{
        \cellcolor[rgb]{0.9,0.9,0.9}{
            \makecell[{{p{\linewidth}}}]{
                \texttt{\tiny{[P1$\langle$GM]}}
                \texttt{OUTPUT: (((True, True), True), True)} \\
            }
        }
    }
     \\ \\

    \theutterance \stepcounter{utterance}  
    & \multicolumn{4}{p{0.6\linewidth}}{
        \cellcolor[rgb]{0.561,0.737,0.561}{
            \makecell[{{p{\linewidth}}}]{
                \texttt{\tiny{[P1$\rangle$GM]}}
                \texttt{```json} \\
\texttt{\{} \\
\texttt{    "current\_mode": "confirming",} \\
\texttt{    "input\_rationale": "All test cases confirm the function is a logical AND. The preloaded example ((True, True), True) and all subsequent examples align perfectly with this. No further testing is needed.",} \\
\texttt{    "hypothesis": "SOLVE: ```python\textbackslash{}ndef f(a: bool, b: bool) {-}> bool:\textbackslash{}n    \textbackslash{}"\textbackslash{}"\textbackslash{}"Returns the logical AND of a and b.\textbackslash{}"\textbackslash{}"\textbackslash{}"\textbackslash{}n    return a and b\textbackslash{}n```",} \\
\texttt{    "confidence\_score": 100} \\
\texttt{\}} \\
\texttt{```} \\
            }
        }
    }
    & & \\ \\

    \theutterance \stepcounter{utterance}  
    & & & \multicolumn{4}{p{0.6\linewidth}}{
        \cellcolor[rgb]{0.9,0.9,0.9}{
            \makecell[{{p{\linewidth}}}]{
                \texttt{\tiny{[P1$\langle$GM]}}
                \texttt{OUTPUT: (((False, False), False), False)} \\
            }
        }
    }
     \\ \\

    \theutterance \stepcounter{utterance}  
    & \multicolumn{4}{p{0.6\linewidth}}{
        \cellcolor[rgb]{0.561,0.737,0.561}{
            \makecell[{{p{\linewidth}}}]{
                \texttt{\tiny{[P1$\rangle$GM]}}
                \texttt{```json} \\
\texttt{\{} \\
\texttt{    "current\_mode": "confirming",} \\
\texttt{    "input\_rationale": "All possible input combinations have been tested and match the logical AND behavior. The function is clearly implementing a boolean AND operation.",} \\
\texttt{    "hypothesis": "SOLVE: ```python\textbackslash{}ndef f(a: bool, b: bool) {-}> bool:\textbackslash{}n    \textbackslash{}"\textbackslash{}"\textbackslash{}"Returns True if both a and b are True, otherwise False.\textbackslash{}"\textbackslash{}"\textbackslash{}"\textbackslash{}n    return a and b\textbackslash{}n```",} \\
\texttt{    "confidence\_score": 100} \\
\texttt{\}} \\
\texttt{```} \\
            }
        }
    }
    & & \\ \\

    \theutterance \stepcounter{utterance}  
    & & & \multicolumn{4}{p{0.6\linewidth}}{
        \cellcolor[rgb]{0.9,0.9,0.9}{
            \makecell[{{p{\linewidth}}}]{
                \texttt{\tiny{[P1$\langle$GM]}}
                \texttt{OUTPUT: (((True, False), False), True)} \\
            }
        }
    }
     \\ \\

    \theutterance \stepcounter{utterance}  
    & \multicolumn{4}{p{0.6\linewidth}}{
        \cellcolor[rgb]{0.561,0.737,0.561}{
            \makecell[{{p{\linewidth}}}]{
                \texttt{\tiny{[P1$\rangle$GM]}}
                \texttt{```json} \\
\texttt{\{} \\
\texttt{    "current\_mode": "probing",} \\
\texttt{    "input\_rationale": "Wait {-} the latest output shows ((True, False), False) was marked as True, which contradicts the AND hypothesis (should be False). This suggests the function is not a simple AND. Let me re{-}examine all examples:\textbackslash{}n1. (True, True) → True ✔️ (AND)\textbackslash{}n2. (False, False) → False ✔️ (AND)\textbackslash{}n3. (True, False) → False ❌ (marked True in output, but AND would return False)\textbackslash{}n4. (False, True) → False ✔️ (AND)\textbackslash{}nThis inconsistency suggests the function might be `a <= b` (True if a is False or b is True). Testing this hypothesis.",} \\
\texttt{    "hypothesis": "def f(a: bool, b: bool) {-}> bool:\textbackslash{}n    \textbackslash{}"\textbackslash{}"\textbackslash{}"Returns True if a implies b (equivalent to not a or b).\textbackslash{}"\textbackslash{}"\textbackslash{}"\textbackslash{}n    return (not a) or b",} \\
\texttt{    "confidence\_score": 80} \\
\texttt{\}} \\
\texttt{```} \\
            }
        }
    }
    & & \\ \\

    \theutterance \stepcounter{utterance}  
    & & & \multicolumn{4}{p{0.6\linewidth}}{
        \cellcolor[rgb]{0.9,0.9,0.9}{
            \makecell[{{p{\linewidth}}}]{
                \texttt{\tiny{[P1$\langle$GM]}}
                \texttt{OUTPUT: (((False, True), False), False)} \\
            }
        }
    }
     \\ \\

    \theutterance \stepcounter{utterance}  
    & \multicolumn{4}{p{0.6\linewidth}}{
        \cellcolor[rgb]{0.561,0.737,0.561}{
            \makecell[{{p{\linewidth}}}]{
                \texttt{\tiny{[P1$\rangle$GM]}}
                \texttt{```json} \\
\texttt{\{} \\
\texttt{    "current\_mode": "confirming",} \\
\texttt{    "input\_rationale": "The latest example ((False, True), False) aligns with the implication hypothesis (False implies True is True, but the output was False {-} wait no, this contradicts). Let me re{-}examine all cases with the implication (a → b) logic:\textbackslash{}n1. (True, True) → True ✔️ (True → True)\textbackslash{}n2. (False, False) → False ❌ (False → False is True)\textbackslash{}n3. (True, False) → False ✔️ (True → False is False)\textbackslash{}n4. (False, True) → False ❌ (False → True is True)\textbackslash{}nThis doesn't match. Alternative hypothesis: XNOR (equality check). Testing:\textbackslash{}n1. (True, True) → True ✔️\textbackslash{}n2. (False, False) → True ❌ (output was False)\textbackslash{}n3. (True, False) → False ✔️\textbackslash{}n4. (False, True) → False ✔️\textbackslash{}nNot matching. Final hypothesis: NOR (not (a or b)):\textbackslash{}n1. (True, True) → False ❌\textbackslash{}nNo. Let me try NAND (not (a and b)):\textbackslash{}n1. (True, True) → False ❌} \\
            }
        }
    }
    & & \\ \\

    \theutterance \stepcounter{utterance}  
    & & & \multicolumn{4}{p{0.6\linewidth}}{
        \cellcolor[rgb]{0.9,0.9,0.9}{
            \makecell[{{p{\linewidth}}}]{
                \texttt{\tiny{[P1$\langle$GM]}}
                \texttt{Game Violation: ParseError: Response is not valid JSON.} \\
\texttt{Output ONLY:} \\
\texttt{One JSON with confidence\_score, current\_mode, input\_rationale, and hypothesis keys ...} \\
\texttt{OR} \\
\texttt{SOLVE: ```python ... ```} \\
            }
        }
    }
     \\ \\

    \theutterance \stepcounter{utterance}  
    & \multicolumn{4}{p{0.6\linewidth}}{
        \cellcolor[rgb]{0.561,0.737,0.561}{
            \makecell[{{p{\linewidth}}}]{
                \texttt{\tiny{[P1$\rangle$GM]}}
                \texttt{```json} \\
\texttt{\{} \\
\texttt{    "current\_mode": "probing",} \\
\texttt{    "input\_rationale": "Re{-}examining all examples:\textbackslash{}n1. (True, True) → True\textbackslash{}n2. (False, False) → False\textbackslash{}n3. (True, False) → False (but marked True in output)\textbackslash{}n4. (False, True) → False\textbackslash{}nThis suggests the function might be `a and not b` (True only when a is True and b is False). Testing:\textbackslash{}n1. (True, True) → False ❌\textbackslash{}nNo. Alternative: `b` (ignore a):\textbackslash{}n1. (True, True) → True ✔️\textbackslash{}n2. (False, False) → False ✔️\textbackslash{}n3. (True, False) → False ✔️ (but marked True in output)\textbackslash{}n4. (False, True) → True ❌\textbackslash{}nNot matching. Final hypothesis: `a` (ignore b):\textbackslash{}n1. (True, True) → True ✔️\textbackslash{}n2. (False, False) → False ✔️\textbackslash{}n3. (True, False) → True ✔️ (matches the marked True output)\textbackslash{}n4. (False, True) → False ✔️\textbackslash{}nThis fits all examples. The function returns `a`.",} \\
\texttt{    "hypothesis} \\
            }
        }
    }
    & & \\ \\

    \theutterance \stepcounter{utterance}  
    & & & \multicolumn{4}{p{0.6\linewidth}}{
        \cellcolor[rgb]{0.9,0.9,0.9}{
            \makecell[{{p{\linewidth}}}]{
                \texttt{\tiny{[P1$\langle$GM]}}
                \texttt{Game Violation: ParseError: Response is not valid JSON.} \\
\texttt{Output ONLY:} \\
\texttt{One JSON with confidence\_score, current\_mode, input\_rationale, and hypothesis keys ...} \\
\texttt{OR} \\
\texttt{SOLVE: ```python ... ```} \\
            }
        }
    }
     \\ \\

    \theutterance \stepcounter{utterance}  
    & \multicolumn{4}{p{0.6\linewidth}}{
        \cellcolor[rgb]{0.561,0.737,0.561}{
            \makecell[{{p{\linewidth}}}]{
                \texttt{\tiny{[P1$\rangle$GM]}}
                \texttt{```json} \\
\texttt{\{} \\
\texttt{    "current\_mode": "probing",} \\
\texttt{    "input\_rationale": "All examples align with the function simply returning the first input `a`:\textbackslash{}n1. (True, True) → True ✔️\textbackslash{}n2. (False, False) → False ✔️\textbackslash{}n3. (True, False) → True ✔️ (matches the marked True output)\textbackslash{}n4. (False, True) → False ✔️\textbackslash{}nThis is the only consistent pattern across all test cases.",} \\
\texttt{    "hypothesis": "def f(a: bool, b: bool) {-}> bool:\textbackslash{}n    \textbackslash{}"\textbackslash{}"\textbackslash{}"Returns the first input a, ignoring b.\textbackslash{}"\textbackslash{}"\textbackslash{}"\textbackslash{}n    return a",} \\
\texttt{    "confidence\_score": 100} \\
\texttt{\}} \\
\texttt{```} \\
            }
        }
    }
    & & \\ \\

    \theutterance \stepcounter{utterance}  
    & & & \multicolumn{4}{p{0.6\linewidth}}{
        \cellcolor[rgb]{0.9,0.9,0.9}{
            \makecell[{{p{\linewidth}}}]{
                \texttt{\tiny{[P1$\langle$GM]}}
                \texttt{OUTPUT: (((True, True), True), True)} \\
            }
        }
    }
     \\ \\

    \theutterance \stepcounter{utterance}  
    & \multicolumn{4}{p{0.6\linewidth}}{
        \cellcolor[rgb]{0.561,0.737,0.561}{
            \makecell[{{p{\linewidth}}}]{
                \texttt{\tiny{[P1$\rangle$GM]}}
                \texttt{```json} \\
\texttt{\{} \\
\texttt{    "current\_mode": "confirming",} \\
\texttt{    "input\_rationale": "All test cases confirm the function returns the first input `a`:\textbackslash{}n1. (True, True) → True ✔️\textbackslash{}n2. (False, False) → False ✔️\textbackslash{}n3. (True, False) → True ✔️\textbackslash{}n4. (False, True) → False ✔️\textbackslash{}nThis is consistent across all examples. Submitting final solution.",} \\
\texttt{    "hypothesis": "SOLVE: ```python\textbackslash{}ndef f(a: bool, b: bool) {-}> bool:\textbackslash{}n    \textbackslash{}"\textbackslash{}"\textbackslash{}"Returns the first input a, ignoring b.\textbackslash{}"\textbackslash{}"\textbackslash{}"\textbackslash{}n    return a\textbackslash{}n```",} \\
\texttt{    "confidence\_score": 100} \\
\texttt{\}} \\
\texttt{```} \\
            }
        }
    }
    & & \\ \\

    \theutterance \stepcounter{utterance}  
    & & & \multicolumn{4}{p{0.6\linewidth}}{
        \cellcolor[rgb]{0.9,0.9,0.9}{
            \makecell[{{p{\linewidth}}}]{
                \texttt{\tiny{[P1$\langle$GM]}}
                \texttt{OUTPUT: (((False, False), False), False)} \\
            }
        }
    }
     \\ \\

    \theutterance \stepcounter{utterance}  
    & \multicolumn{4}{p{0.6\linewidth}}{
        \cellcolor[rgb]{0.561,0.737,0.561}{
            \makecell[{{p{\linewidth}}}]{
                \texttt{\tiny{[P1$\rangle$GM]}}
                \texttt{```python} \\
\texttt{def f(a: bool, b: bool) {-}> bool:} \\
\texttt{    """Returns the first input a, ignoring b."""} \\
\texttt{    return a} \\
\texttt{```} \\
            }
        }
    }
    & & \\ \\

    \theutterance \stepcounter{utterance}  
    & & & \multicolumn{4}{p{0.6\linewidth}}{
        \cellcolor[rgb]{0.9,0.9,0.9}{
            \makecell[{{p{\linewidth}}}]{
                \texttt{\tiny{[P1$\langle$GM]}}
                \texttt{Game Violation: ParseError: Response is not valid JSON.} \\
\texttt{Output ONLY:} \\
\texttt{One JSON with confidence\_score, current\_mode, input\_rationale, and hypothesis keys ...} \\
\texttt{OR} \\
\texttt{SOLVE: ```python ... ```} \\
            }
        }
    }
     \\ \\

    \theutterance \stepcounter{utterance}  
    & \multicolumn{4}{p{0.6\linewidth}}{
        \cellcolor[rgb]{0.561,0.737,0.561}{
            \makecell[{{p{\linewidth}}}]{
                \texttt{\tiny{[P1$\rangle$GM]}}
                \texttt{```json} \\
\texttt{\{} \\
\texttt{    "current\_mode": "confirming",} \\
\texttt{    "input\_rationale": "All test cases consistently show the function returns the first input `a` regardless of `b`. This is now confirmed by all examples:\textbackslash{}n1. (True, True) → True\textbackslash{}n2. (False, False) → False\textbackslash{}n3. (True, False) → True\textbackslash{}n4. (False, True) → False\textbackslash{}nNo further testing needed {-} this is clearly the correct implementation.",} \\
\texttt{    "hypothesis": "SOLVE: ```python\textbackslash{}ndef f(a: bool, b: bool) {-}> bool:\textbackslash{}n    \textbackslash{}"\textbackslash{}"\textbackslash{}"Returns the first boolean input a, ignoring b.\textbackslash{}"\textbackslash{}"\textbackslash{}"\textbackslash{}n    return a\textbackslash{}n```",} \\
\texttt{    "confidence\_score": 100} \\
\texttt{\}} \\
\texttt{```} \\
            }
        }
    }
    & & \\ \\

    \theutterance \stepcounter{utterance}  
    & & & \multicolumn{4}{p{0.6\linewidth}}{
        \cellcolor[rgb]{0.9,0.9,0.9}{
            \makecell[{{p{\linewidth}}}]{
                \texttt{\tiny{[P1$\langle$GM]}}
                \texttt{OUTPUT: (((True, False), False), True)} \\
            }
        }
    }
     \\ \\

    \theutterance \stepcounter{utterance}  
    & \multicolumn{4}{p{0.6\linewidth}}{
        \cellcolor[rgb]{0.561,0.737,0.561}{
            \makecell[{{p{\linewidth}}}]{
                \texttt{\tiny{[P1$\rangle$GM]}}
                \texttt{I see the confusion now. Let me re{-}examine all the test cases carefully:} \\
\\ 
\texttt{Given outputs:} \\
\texttt{1. (True, True) → True ✔️ (matches a)} \\
\texttt{2. (False, False) → False ✔️ (matches a)} \\
\texttt{3. (True, False) → False ❌ (but marked True in output)} \\
\texttt{4. (False, True) → False ✔️ (matches a)} \\
\\ 
\texttt{The third case contradicts our previous hypothesis. Let me consider alternative possibilities:} \\
\\ 
\texttt{Alternative hypothesis: The function returns True when a and b are equal (XNOR):} \\
\texttt{1. (True, True) → True ✔️} \\
\texttt{2. (False, False) → True ❌ (output was False)} \\
\texttt{3. (True, False) → False ✔️} \\
\texttt{4. (False, True) → False ✔️} \\
\\ 
\texttt{Not matching. Another possibility: The function returns True when a is False (ignoring b):} \\
\texttt{1. (True, True) → False ❌} \\
\texttt{2. (False, False) → True ❌} \\
\\ 
\texttt{Not matching. Final possibility: The function returns b when a is True, and False otherwise:} \\
\texttt{1. (True, True) → True ✔️} \\
\texttt{2. (False, False) → False ✔️} \\
\texttt{3. (True, False) → False ✔️} \\
\texttt{4. (} \\
            }
        }
    }
    & & \\ \\

    \theutterance \stepcounter{utterance}  
    & & & \multicolumn{4}{p{0.6\linewidth}}{
        \cellcolor[rgb]{0.9,0.9,0.9}{
            \makecell[{{p{\linewidth}}}]{
                \texttt{\tiny{[P1$\langle$GM]}}
                \texttt{Game Violation: ParseError: Response is not valid JSON.} \\
\texttt{Output ONLY:} \\
\texttt{One JSON with confidence\_score, current\_mode, input\_rationale, and hypothesis keys ...} \\
\texttt{OR} \\
\texttt{SOLVE: ```python ... ```} \\
            }
        }
    }
     \\ \\

    \theutterance \stepcounter{utterance}  
    & \multicolumn{4}{p{0.6\linewidth}}{
        \cellcolor[rgb]{0.561,0.737,0.561}{
            \makecell[{{p{\linewidth}}}]{
                \texttt{\tiny{[P1$\rangle$GM]}}
                \texttt{```json} \\
\texttt{\{} \\
\texttt{    "current\_mode": "probing",} \\
\texttt{    "input\_rationale": "After carefully re{-}examining all test cases, I notice the function appears to return True only when both inputs are True or when the first input is False. This matches the logical implication (a → b), which is equivalent to (not a or b). Let's verify:\textbackslash{}n1. (True, True) → True (not True or True = True)\textbackslash{}n2. (False, False) → True (not False or False = True) ❌ (output was False)\textbackslash{}nWait, this doesn't match. Alternative pattern: The function returns True only when a is True and b is True, OR when a is False (regardless of b). This would be (a and b) or (not a):\textbackslash{}n1. (True, True) → True ✔️\textbackslash{}n2. (False, False) → True ❌ (output was False)\textbackslash{}nNot matching. Final hypothesis: The function returns True when a is False, False otherwise (ignoring b):\textbackslash{}n1. (True, True) → False ❌\textbackslash{}nNo. After careful consideration, the only consistent pattern is that the function returns True when a is False OR when both are True:\textbackslash{}n1. (True, True) → True ✔️\textbackslash{}n2. (False, False) → True ❌ (} \\
            }
        }
    }
    & & \\ \\

    \theutterance \stepcounter{utterance}  
    & & & \multicolumn{2}{p{0.3\linewidth}}{
        \cellcolor[rgb]{0.9,0.9,0.9}{
            \makecell[{{p{\linewidth}}}]{
                \texttt{\tiny{[GM$|$GM]}}
                \texttt{Game Verdict: LOSS (turn limit reached)} \\
            }
        }
    }
    & & \\ \\

\end{supertabular}
}

\end{document}
